\documentclass{ximera}

\begin{document}
    \title{Limit Practice Problems}
    \begin{abstract} Time to practice! \end{abstract}
    \maketitle

\section*{Additional Practice Problems}

Try solving the problems below for practice, without a calculator.  \textbf{Enter exact answers only}.  To enter an answer
containing a radical, use the \textbf{sqrt( )} function.  For example, type sqrt(2) for $\sqrt{2}.$  To enter an answer containing 
$\infty,$ type \textbf{infinity} or \textbf{oo} (two lowercase letter o's).

\begin{problem}
Evaluate the limit.

\[ \lim_{x \rightarrow 2} \sin \left( \frac{\pi x}{6} \right)  \begin{prompt} = \answer{\sqrt{3}/2} \end{prompt} \]

%To accept a decimal approximation or exact answer, use a tolerance of 3 decimal places for trigonometric functions: \answer[tolerance=0.001]{\sqrt{3}/2}

Does the limit exist?

\begin{multipleChoice}
    \choice[correct]{Yes.}
    \choice{No.}
  \end{multipleChoice}

\end{problem}


\begin{problem}
Evaluate the limit.

\[ \lim_{x\to -3} \frac{x+3}{x^2-9} =\answer{-1/6} \]


Does the limit exist?


\begin{multipleChoice}
    \choice[correct]{Yes.}
    \choice{No.}
  \end{multipleChoice}

\end{problem}


\begin{problem}
Evaluate the limit.

\[\lim_{x\to 10^+} \frac{3+x}{10-x}  \begin{prompt}=\answer{-\infty}\end{prompt}\]


Does the limit exist?

\begin{multipleChoice}
    \choice{Yes.}
    \choice[correct]{No.}
  \end{multipleChoice}

\end{problem}


\begin{problem}
Evaluate the limit.

\[ \lim_{\theta\to 0} \frac{2 \sin \theta (1-\cos \theta)}{5 \theta^2} \begin{prompt}=\answer{0}\end{prompt}\]


Does the limit exist?

\begin{multipleChoice}
    \choice[correct]{Yes.}
    \choice{No.}
  \end{multipleChoice}

\begin{hint}
\textbf{Use the properties of limits to rewrite the limit as the product of limits that you can evaluate.  See the 
Evaluating Limits Analytically section on the previous page.}
\end{hint}


\end{problem}

\begin{problem}
Evaluate the limit.

\[ \lim_{t\to 0^+} \frac{1}{t^2} \begin{prompt}=\answer{\infty}\end{prompt}\]


Does the limit exist?

\begin{multipleChoice}
    \choice{Yes.}
    \choice[correct]{No.}
  \end{multipleChoice}


\end{problem}


\begin{problem}
Evaluate the limit.

\[ \lim_{x\to 7} \frac{x-7}{\sqrt{x} - \sqrt{7}} \begin{prompt}=\answer{2\sqrt{7}}\end{prompt}\]


Does the limit exist?
\begin{multipleChoice}
    \choice[correct]{Yes.}
    \choice{No.}
  \end{multipleChoice}

\begin{hint}
\textbf{Rationalize the denominator.}
\end{hint}

\end{problem}



\begin{problem}
During a rapid intravenous infusion, the concentration of a therapeutic drug in a patient's liver tissue during the first hour 
is modeled by the function

\[ C(t) = \frac{3}{1 - t} \]

where $t$ is the time in hours ($0 \leq t < 1$) and $C(t)$ is measured in milligrams per gram (mg/g).

As time approaches $1$ hour from the left, find the limit of the drug concentration.

\[ \lim_{t \to 1^-} C(t) = \answer{\infty} \]

Based on this limit, what is happening to the drug concentration inside the tissue as time approaches $1$ hour?

\begin{multipleChoice}
  \choice{The drug concentration drops to zero.}
  \choice{The drug concentration stabilizes at a safe, constant level.}
  \choice[correct]{The drug concentration spikes rapidly toward dangerously high levels.}
\end{multipleChoice}
\end{problem}

\end{document}